\section{Introduction}\label{sec:intro}

An attractor is called hidden if its basin of attraction does not intersect small neighborhoods of equilibria; otherwise, the attractor is self-excited~\cite{LeonovKuznetsov2013,Leonov2011PLA,Kuznetsov2020JCSSI}. Hidden oscillations occur in Chua circuits~\cite{Leonov2012PhysD,Stankevich2017,Kuznetsov2023Chua}, in control systems with saturation~\cite{Andrievsky2017Aircraft}, in the Aizerman and Kalman problems~\cite{Bragin2011}, in generalized Lorenz and Rabinovich systems~\cite{Chen2017OnePath,Kuznetsov2018Rabinovich,Kuznetsov2020Lorenz}, in systems without equilibria~\cite{JafariSprottGolpayegani2013,Wei2011}, with a single stable equilibrium~\cite{WangChen2012,Molaie2013}, with lines and surfaces of equilibria~\cite{JafariSprott2013,Jafari2016Surfaces}, in two-dimensional maps~\cite{Jiang2016Maps,WangDing2018Map,Huynh2019}, and in neural network models~\cite{DancaKuznetsov2017Hopfield}; see the reviews~\cite{Dudkowski2016,JafariSprottNazarimehr2015,GuanXie2025}. A hidden attractor usually coexists with other attractors, so searching for it is a problem of multistable dynamics~\cite{Feudel2008,PisarchikFeudel2014}. Attractors whose basins occupy a small fraction of the state space are called rare~\cite{Chudzik2011Rare,Yue2021RareDuffing}.

The difficulty stems from the definition. The standard procedure for finding attractors launches trajectories from neighborhoods of unstable equilibria along their unstable manifolds, and by definition it cannot find a hidden attractor. Analytical--numerical localization uses harmonic balance and the describing function method, followed by continuation in a parameter~\cite{Leonov2011PLA,KuznetsovLeonovVagaitsev2010,Kuznetsov2017IFAC}. Geometric heuristics propose starting sets: perpetual points~\cite{Prasad2015,Dudkowski2015PP,Dudkowski2017PP}, connecting curves~\cite{Gilmore2010CC,GuanXie2021CC,Naseri2021} and surfaces~\cite{Song2023CS}, coordinate-wise zero-velocity sets~\cite{Godara2018CVS}, and parameter switching~\cite{Danca2018Switching}; evolutionary search over initial conditions is used as well~\cite{Zelinka2016,Zelinka2021Chapter}. The general-purpose tool is the estimation of basins by sampling initial conditions~\cite{Menck2013,DatserisWagemakers2022,Datseris2023Framework,StenderHoffmann2022,Gelbrecht2020MCBB}, complemented by set-oriented methods~\cite{Hsu1987,DellnitzHohmann1997}. A sample of $n$ points detects an attractor with basin fraction $p$ with probability $1-(1-p)^n$, so rare attractors are found either with very large samples or by continuation from parameter values where they are not rare~\cite{Brzeski2017,Datseris2023Framework,Datseris2026Global,Haerter2026BBT}. The heuristics have known limitations: perpetual points may not exist or may lead elsewhere~\cite{Nazarimehr2017PP,Jafari2015LimitPP,Jafari2026NoPP}, and connecting curves need not pass through the basin of the hidden attractor~\cite{GuanXie2025}.

Assessments of localization methods are almost always qualitative: the attractor was found from a certain point. How many integrations a method needs, with what probability it finds the target within a given budget, and whether it beats uniform sampling at all are rarely reported~\cite{GuanXie2025,Zelinka2016}. Hiddenness and rarity are, moreover, different properties: in well-known examples, the basin of a hidden attractor often occupies most of the natural region of initial conditions (\cref{sec:bench}), and what makes a search hard is a small basin measure $p$, not the absence of a connection to equilibria. Long-lived transient chaos is also sometimes declared a hidden attractor~\cite{Munmuangsaen2018,SprottMunmuangsaen2018Comment,Danca2016Transient,YorkeYorke1979,GrebogiOttYorke1983,LaiTel2011}, so a comparison of methods needs an explicit verification protocol.

If the target basin is small, finding the attractor is the detection of a rare event in the space of initial conditions, a setting with well-developed tools: multilevel splitting~\cite{Glasserman1999,CerouGuyader2007,Cerou2019AMS} and its static counterpart, subset simulation~\cite{AuBeck2001}, the cross-entropy method~\cite{Rubinstein1999CE}, sequential design with active learning of the failure domain~\cite{Bect2012,Echard2011AKMCS}, and novelty search~\cite{LehmanStanley2011,MouretClune2015,Pugh2016QD,CullyDemiris2018}. Such algorithms have been applied to extreme events in climate models~\cite{Ragone2018,Lestang2018TAMS} and to high-dimensional dynamical systems~\cite{Baars2021AMS}. Machine-learning methods estimate the fractions of already known basins~\cite{Wang2022HAL,Cao2025AL,Yang2024PRR,DNNBasins2021,Gameiro2025,Haerter2026BBT}. To the best of our knowledge, these tools have not been applied systematically to the search for hidden attractors, and no bound is known on how much any search method can gain over uniform sampling.

This paper proposes a search method for rare hidden attractors whose cost never exceeds that of uniform sampling by more than a fixed factor, and answers the question of when a hidden attractor can be found faster than by brute force, and by how much. Its contributions are as follows.
\begin{enumerate}
\item \emph{Information bound} (\cref{sec:problem}). We measure the cost of a search by the number of integrations needed to detect the target with confidence $1-\delta$ and prove that if everything an algorithm observes, apart from hitting the target, is constant on a set $Q$ containing the target basin, then no algorithm can guarantee a gain over uniform sampling larger than about $\ln(1/\delta)/((1-\delta)\mu(Q))$, which is $2.6/\mu(Q)$ for $\delta=0.1$ (\cref{thm:info}). Without information beyond the target indicator, adaptivity is worth at most a factor of 2.6; escape times are worth at most $2.6/\mu(Q_\infty)$, where $Q_\infty$ is the nonescaping set (\cref{cor:channels}).
\item \emph{Theory of the gain along a reaction coordinate} (\cref{sec:theory}). A cost model of multilevel splitting classifies coordinates into those that give an order-of-growth gain, a constant-factor gain, or a loss (\cref{prop:cost}). For $k$ uncoupled units, ideal splitting along the escape time attains the information bound up to the factor 2.6, and both the bound and the gain grow geometrically in $k$ with the base $1/q$, where $q$ is the measure of the nonescaping set of a single unit, which is measurable before any search (\cref{prop:product}). The bound is thus sharp up to a constant, and rare hidden regimes of systems of units are the class where adaptive search can gain by orders of magnitude.
\item \emph{A reaction coordinate and a portfolio with a worst-case guarantee} (\cref{sec:method}). The transient duration ranks escaping trajectories by their escape time and bounded ones by how slowly they settle onto their regime, so it also carries information where escape times are constant. Splitting along it, active learning of basin boundaries, and a uniform share $\theta$ are combined so that the cost is never more than $1/\theta$ times that of uniform sampling, up to an additive constant, and the scaling exponent of every autonomous arm is inherited (\cref{prop:portfolio}). In particular, stagger-and-step chains on the same coordinate become a safe arm (\cref{cor:sas}).
\item \emph{Benchmark and verification protocol} (\cref{sec:bench}). The benchmark consists of three-dimensional flows, maps, four-dimensional systems, and synthetic families with analytically known basin measure. It comes with a protocol that separates attractors from transient chaos and from integration-step artifacts, and with a hiddenness measure $h_\varepsilon$ that grades hidden attractors by the scale at which the equilibria cease to lead to them.
\item \emph{Experiments that test the theory} (\cref{sec:results}). On coupled copies of a system with a hidden attractor, the gain grows geometrically, from 65 to $1.3\times10^4$ for two to four copies, with base 14 compared with the base 21 predicted for ideal splitting from a single copy, and stays below the information bound. Where all trajectories are bounded and escape times carry no information, the settling part of the coordinate gives a gain of 7.4 on a Chua circuit and gains of 55--360 on systems of two and three Chua circuits, and the gain carries over to wall time. On a synthetic family with a transient signal on all scales, the cost grows as $p^{-0.45}$ instead of $p^{-1}$; on the rare hidden attractors of the benchmark that have a transient signal, the gain is 3--7, whereas evolutionary search is no better than uniform sampling. The failures are those predicted by the cost model. The enumeration of regimes with the protocol reveals a hidden chaotic attractor of two coupled Chua circuits that a single circuit does not have.
\end{enumerate}
The equations of the benchmark systems, verification episodes, ablations, enlarged regions, the discretization limit, channel diagnostics, and implementation details are given in the Supplementary Material (Sections~\ref{app:systems}--\ref{app:implementation}).
